\documentclass[12pt,a4paper]{article}

\usepackage[utf8]{inputenc}
\usepackage[T1]{fontenc}
\usepackage[margin=2.5cm]{geometry}
\usepackage{hyperref}
\usepackage{microtype}
\usepackage{parskip}
\usepackage{lmodern}
\usepackage{amsmath}

\hypersetup{colorlinks=true, urlcolor=blue!60!black}

\pagestyle{empty}

\begin{document}

\begin{flushright}
Kundai Farai Sachikonye \\
Auf der Lichtung 19, 82178 Puchheim, Germany \\
\href{https://github.com/fullscreen-triangle}{github.com/fullscreen-triangle} \\
\texttt{kundai.sachikonye@bitspark.com} \\
(+49) 1626386344 \\[6pt]
18 May 2026
\end{flushright}

\vspace{1em}

Prof.\ Dr.\ Andre Franke \\
Institut für Klinische Molekularbiologie (IKMB) \\
Universitätsklinikum Schleswig-Holstein, Campus Kiel \\
Ref.\ 28608

\vspace{1.5em}

\textbf{Re: PhD Position --- Computational Immunology / IBD Genomics (Ref.\ 28608)}

\vspace{1em}

Dear Prof.\ Franke,

I am applying for the PhD position in computational immunology at IKMB (Ref.\ 28608).
The position covers immune repertoire analysis, single-cell sequencing, and the
microbiome--host genetics interface in intestinal inflammation --- precisely the
application domain for a computational immunopeptidology toolkit I have developed
and deployed over the past four years.
The specific analytical questions the position describes --- which epitopes drive
mucosal T cell responses, how immune repertoires shift across IBD patient cohorts,
how host genetic variation modulates inflammatory microbiome interactions ---
are computationally tractable with the spectral methods and Bayesian inference
infrastructure I have built.

\textbf{Computational immunopeptidology for IBD T cell biology.}
The \href{https://syndrome-seven.vercel.app/tools}{Syndrome} framework applies
3-channel discrete Fourier transforms (hydropathy, volume, charge) to compute
spectral complementarity $\rho$ between any two molecular sequences in
$\mathcal{O}(cL\log L)$.
For IBD, the three directly applicable tools are:
\begin{itemize}
  \item \textbf{MHC binding predictor}: $\rho \geq 0.72$ for binding prediction;
    AUC\,$=0.81$ vs.\ NetMHCpan\,$=0.68$ on 2,847 IEDB peptides
    ($r=0.73$, $p<10^{-12}$).
    Applicable to patient-specific HLA scanning: given a patient's HLA haplotype
    and the microbial peptide landscape from their metagenomics data, predict
    which microbial epitopes will be presented and drive mucosal T cell activation.
  \item \textbf{Neoantigen / cross-reactive epitope identification}: dual threshold
    $\Delta\rho_\mathrm{self} > 0.15$ (sufficient divergence from self) and
    $\rho_\mathrm{MHC} > 0.72$ (sufficient binding affinity).
    In IBD, this identifies microbial epitopes with high cross-reactivity risk
    to intestinal self-antigens --- the computational screen for molecular
    mimicry candidates.
  \item \textbf{Escape predictor}: systematic single amino acid substitution scan
    scoring $E(v) = -\max\rho(v, \mathrm{self}) + w \cdot \rho(v, \mathrm{target})$.
    For immune repertoire analysis, this maps how observed TCR sequences relate to
    the epitope landscape: which receptor variants maintain target recognition
    while minimising autoreactive cross-reactivity.
\end{itemize}
These tools are live at \href{https://syndrome-seven.vercel.app/tools}{syndrome-seven.vercel.app}
and require no local installation, making them directly usable by wet-lab
collaborators without computational infrastructure.

\textbf{IBD genomics and microbiome--host genetics.}
The \href{https://github.com/fullscreen-triangle/gospel}{Gospel} framework
implements large-scale genomics analysis at $10^6$ variants per second:
somatic and germline variant calling, RNA-seq quantification and differential
expression, Bayesian network inference for multi-omics integration, and
Gaussian mixture models for population structure.
Validated against 1000~Genomes Phase~3, gnomAD~v3.1.2, and UK~Biobank ---
the same cohort scales as IBD GWAS studies.
The Bayesian network module performs multi-omics integration directly:
host genotype variation as upstream nodes conditioning metagenomics and
transcriptomics distributions downstream, with posterior inference across
the joint graph.
This is the computational structure for microbiome--host genetics analysis:
identifying which host genetic variants (IBD GWAS loci) modulate microbial
community composition and mucosal inflammatory gene expression simultaneously,
rather than correlating the two separately.

\textbf{Disease state trajectory and therapeutic target selection.}
The Syndrome framework's disease state layer represents patient state as a vector
$\mathbf{D} \in [0,1]^8$ with dimensions including immune activation ($D_\mathrm{Imm}$),
mucosal barrier integrity ($D_\mathrm{Bar}$), ATP synthesis coherence ($D_\mathrm{ATP}$),
and circadian regulatory coherence ($D_\mathrm{Circ}$).
Trajectory inference uses a Bayesian $P(R \mid Q)$ posterior gating structure,
giving a principled model of disease progression rather than a cross-sectional
score.
For minimum-intervention target identification --- which single pathway
perturbation most efficiently restores the disease state vector toward
remission --- the framework uses a sparse $\ell_1$ linear program, validated
at AUC\,$=1.00$ across 25/25 validation experiments.
In the IBD context, this provides a computational framework for
personalised therapeutic target prioritisation from single-cell multi-omics
data: patient state characterisation from scRNA-seq and TCR-seq jointly,
followed by LP-based identification of the minimal intervention set.

\textbf{Single-cell methods and repertoire analysis.}
My primary single-cell experience is through the Bayesian inference layer:
variational Bayes posterior approximation at $\mathcal{O}(10^6)$ latent
variable scale (Gospel), probabilistic graphical models for multi-omics
integration, and Gaussian processes for trajectory inference.
Scanpy and Seurat are familiar from MSc training and independent study;
the statistical models underlying UMAP-based clustering and RNA velocity
are directly related to the Bayesian dimensionality reduction infrastructure
I have built.
TCR/BCR clonotype analysis --- the combinatorial structure of V(D)J recombination,
clonal expansion statistics, and cross-patient repertoire comparison --- is a
sequence analysis problem where the spectral complementarity methods in Syndrome
apply directly: $\rho$ as a distance measure between TCR sequences, enabling
clustering of functionally similar clonotypes without requiring MSA.
The PhD provides the experimental immunology context (patient samples, immune
phenotyping, metagenomics protocols) that gives these computational methods
their biological grounding.

\textbf{Background.}
I hold an MSc in Computational Biology (Universität Tübingen, 2018--2019)
and conducted doctoral research in Computational Lipidomics at TUM\slash
Bayerisches Forschungsinstitut (2019--2021), working on ML for
large-scale MS omics data in clinical cohorts.
I am legally resident in Germany with full work authorisation.
English is native; German is conversational.

\vspace{1.5em}
Yours sincerely,

\vspace{2.5em}
\textbf{Kundai Farai Sachikonye}

\end{document}
